% from main.tex
%%%%%%%%%%  scan idx 530  |  folio 517  %%%%%%%%%%
Or d'après VI 4.8, $K^\circ(X)$ est la $K^\circ(S)$-algèbre
engendrée par les $\xi_i-1$, soumis aux relations déduites de :
\[
\gamma_t(E-n)=\prod_{i=1}^n\gamma_t(\xi_i-1).
\]
Donc $K^\circ(X)/I=K^\circ(S)/J$, où $J$ est l'idéal engendré par les $\gamma^j(E-n)$ pour $j=1,\ldots,n$, donc le $\lambda$-idéal engendré par $E-n$. Par suite,
\[
K_\bullet(Y)\simeq K_\bullet(S)/J K_\bullet(S).
\]
On achève la démonstration par le lemme :

\textbf{Lemme 4.2.} Soient $S$ un schéma noethérien, $B$ un sous-groupe parabolique d'un $S$-groupe réductif, $P$ un torseur sous le radical unipotent $U$ de $B$; alors on a un isomorphisme :
\[
K_\bullet(S)\xrightarrow{\sim}K_\bullet(P).
\]
En effet, il existe une suite de sous-schémas en groupes de $U$ :
\[
U_0=U\supset U_1\supset U_2\supset\cdots\supset U_p=e,
\]
fermés dans $B$, finie car $B$ est noethérien, et tels que, pour tout $i$, $U_i/U_{i+1}$ soit isomorphe à un $S$-groupe vectoriel (SGA 3 XXVI 2.1). Par suite, $P\to S$ se factorise par une suite de torseurs sous des groupes vectoriels, d'où le résultat d'après 3.4.

\textbf{Corollaire 4.3.} Soient $S$ un schéma noethérien, $G=\mathrm{GL}(n)_S$. Alors l'homomorphisme canonique $K_\bullet(S)\longrightarrow K_\bullet(G)$ est un isomorphisme.

Résulte immédiatement de 4.1, appliqué à $P=G$, puisqu'alors $E=\mathcal O_S^n$.

